
\section{Homework 7}

\subsection{傅里叶变换性质}

\task{必做题}

(1)$e^{-2t}u(t)$的频谱为$\frac{1}{2+j\omega}$, $\cos^2(t)$的频谱为$\pi\delta(\omega)-
\frac{\pi}{2}[\delta(\omega-2)+\delta(\omega+2)]$, 因此$F(\omega)$为二者的卷积,
$F(\omega)=\pi[\frac{1}{2+j\omega}-\frac{1}{2}(\frac{1}{2+j(\omega-2)}+\frac{1}{2+j(\omega+2)})]$.

(2) $F_1(\omega)=\frac{1}{3}[F(-\frac{\omega}{3})\cdot e^{j(7/3)\omega}]$.

(3) 升余弦信号的Fourier变换为$\frac{E\tau}{2}\frac{Sa(\frac{\omega t}{2})}{1-(\frac{\omega t}{2\pi})^2}$, $\sin(\omega_0 t)$的Fourier变换为$\pi(\delta(\omega-\omega_0)-\delta(\omega+\omega_0))$, 结果为二者卷积:

\begin{equation*}
    \frac{\pi E\tau}{2}\left[\frac{Sa(\frac{(\omega-\omega_0) t}{2})}{1-(\frac{(\omega-\omega_0) t}{2\pi})^2}-\frac{Sa(\frac{(\omega+\omega_0) t}{2})}{1-(\frac{(\omega+\omega_0) t}{2\pi})^2}\right]
\end{equation*}

(4) $f_1(t)=4tf(0.5t-1)+2f(0.5t-1)$. 其中$f(0.5t-1)$的Fourier变换为$2F(2\omega)e^{-2j\omega}$, $tf(0.5t-1)$的Fourier变换为$4e^{-2j\omega}(F(2\omega)+jF'(2\omega))$, 因此结果为$4e^{-2j\omega}[5F(2\omega)=4jF'(2\omega)]$

(5) $f_1$的Fourier变换为$\frac{E\tau}{2}Sa^2(\frac{\omega\tau}{4})$, 则$f$的Fourier变换为$f_1$的Fourier变换的平方, 即$[\frac{E\tau}{2}Sa^2(\frac{\omega\tau}{4})]^2$

(6) $\frac{\sin(\pi t)}{\pi t}$的Fourier变换为$u(t+\pi)-u(t-\pi)$, 则$\frac{\sin(\pi (t-1))}{\pi (t-1)}$的Fourier变换为$[u(t+\pi)-u(t-\pi)]e^{-j\omega}$, 则结果为二者的卷积: 
$$\left\{\begin{array}{ll}
    j(1+e^{-j(\omega+\pi)}) & -2\pi \le \omega < 0 \\
    j(1-e^{-j(\omega-\pi)}) & 0 \le \omega < 2\pi \\
    0 & \text{otherwise}
\end{array}\right.$$

\subsection{傅里叶反变换}

\task{必做题}

(1) 设高为$E$, 底为$\omega_0$的对称三角频谱的反Fourier变换结果为$f_0(t)$, 则最终结果为$f_0(t)[e^{-j\frac{3}{2}\omega_0t}+e^{j\frac{\omega_0}{2}t}]$

由于相同形状的三角频谱的Fourier变换结果为$F_0(\omega)=\frac{E\omega_0}{2}Sa^2(\frac{\omega_0\omega}{4})$, 则$f_0(t)=\frac{1}{2\pi}F_0(-t)=\frac{E\omega_0}{4\pi}Sa^2(\frac{\omega_0\omega}{4})$, 最终结果为$\frac{E\omega_0}{4\pi}Sa^2(\frac{\omega_0\omega}{4})[e^{-j\frac{3}{2}\omega_0t}+e^{j\frac{\omega_0}{2}t}]$

(2) $1/(1+j\omega)$的Fourier反变换为单边指数$e^{-t}u(t)$, 则$F(\omega)$的Fourier反变换结果为其自身卷积, 即$te^{-t}u(t)$

(3) $e^tu(t)$的Fourier变换结果为$1/(1+j\omega)$, 两边做反褶得$F(\omega)$的Fourier反变换为$e^{-t}u(-t)$

\subsection{周期信号的频谱}

\task{必做题}

(1)$f(t)$的Fourier变换为$\displaystyle\frac{\pi}{4}\sum_{-\infty}^{+\infty}Sa(\frac{\pi n}{4})\delta(\omega-2\pi n)+\frac{\pi}{2}\sum_{-\infty}^{+\infty}Sa(\frac{\pi n}{2})\delta(\omega-2\pi n)$

(2) $$\frac{\pi}{T_1}\sum_{-\infty}^{+\infty}[\frac{2E}{\tau}(T_1-\frac{\tau}{2})^2Sa^2(\frac{\pi n}{2T_1}(T_1-\frac{\tau}{2}))-
\frac{2ET_1^2}{\tau}Sa^2(\frac{\pi}{2}nOl+2\pi E\delta(n\frac{\pi}{T_1}))]\delta(\omega-n\frac{\pi}{T_1})$$

\subsection{证明题}

\task{必做题}

对$f$求两阶导数:$\displaystyle\frac{\mathrm{d}^2}{\mathrm{d}t^2}f(t)=\sum_{i=0}^{n}(\frac{x_{i+1}-x_{i}}{t_{i+1}-t_i}-\frac{x_i-x_{i-1}}{t_i-t_{i-1}})\delta(t-t_i)$, 其中规定$x_{-1}=0,t_{-1}<t_0$. 其频谱为$\displaystyle\sum_{i=0}^{n}(\frac{x_{i+1}-x_{i}}{t_{i+1}-t_i}-\frac{x_i-x_{i-1}}{t_i-t_{i-1}})e^{-j\omega t_i}$, 因此原信号的频谱为$\displaystyle\frac{1}{\omega^2}\sum_{i=0}^{n}(-\frac{x_{i+1}-x_{i}}{t_{i+1}-t_i}+\frac{x_i-x_{i-1}}{t_i-t_{i-1}})e^{-j\omega t_i}$, 令$X_i=\displaystyle -\frac{x_{i+1}-x_{i}}{t_{i+1}-t_i}+\frac{x_i-x_{i-1}}{t_i-t_{i-1}}$, 则频谱为$\displaystyle\frac{1}{\omega^2}\sum_{i=0}^{n}X_ie^{-j\omega t_i}$

\subsection{信号的带宽}

由信号带宽的定义, 有:

\begin{align*}
    \frac{1}{2\pi}\int_{-W_{90}}^{W_{90}}\left|\frac{1}{a+j\omega}\right|^2\,\mathrm{d}\omega&=0.9\int_{-\infty}^{+\infty}|f(t)|^2\,\mathrm{d}t \\
    \frac{\arctan(W_{90}/a)}{a\pi}&=0.9\cdot\frac{1}{2a} \\
    W_{90}&=a\tan\left(\frac{9}{20}\pi\right)
\end{align*}
